\section{What is repaired, and what is not}
\label{sec:scope}

\begin{theorem}
  \label{thm:assembled}
  Let \(X\dashrightarrow X'\) be a standard flip with \(r=s+1\) and
  \(s\ge1\).  Then \cite[Theorems~1.2, 1.4, 9.9 and~9.14]{SS} and
  \cite[Corollary~1.5]{SS} hold for \(X\).  In particular they hold for the
  blow-up of a smooth projective variety along a smooth centre of
  codimension two, the case \((r,s)=(2,1)\).
\end{theorem}

\begin{proof}
  Corollary~\ref{cor:presentation} gives \eqref{eq:ss37} and
  \cite[Theorem~1.2]{SS}, and makes \cite[(39)--(40)]{SS} available, so the
  modified extremal \(J\)-function used from \cite[Section~7]{SS} onward is
  the one the argument assumes.  \cite[Proposition~7.5]{SS},
  \cite[Corollary~7.6]{SS} and \cite[Proposition~7.8]{SS} then hold with
  \eqref{eq:ss64} replaced by \eqref{eq:center}, by Lemma~\ref{lem:center};
  \cite[Propositions~8.2 and~8.3]{SS} and \cite[Corollary~8.4]{SS} hold as
  printed, by Lemma~\ref{lem:ambient}.

  The reductions of \cite[Sections~9.1--9.4]{SS} are then unchanged.
  \cite[Proposition~9.1]{SS} uses the presentation \eqref{eq:ss37} and the
  extremal Dubrovin connection; \cite[Proposition~9.2]{SS} combines it with
  \cite[Propositions~7.8 and~8.4]{SS}; \cite[Corollary~9.3]{SS} takes a
  partial nonequivariant limit; \cite[Section~9.2]{SS} removes the splitting
  assumption on \(V,V'\), giving \cite[Corollary~9.7]{SS};
  \cite[Section~9.3]{SS} passes from the local model to \(X\), giving
  \cite[Theorem~9.9]{SS}; and \cite[Section~9.4]{SS} compares the
  Fourier--Mukai transform with the local model, giving
  \cite[Theorem~9.14]{SS}.  None of these steps constrains \(\nu\) beyond
  \(r>s\).  The two halves of \cite[Theorem~1.4]{SS} are
  \cite[Theorems~9.9 and~9.14]{SS}, and \cite[Corollary~1.5]{SS} is the
  translation of \cite[Theorem~1.4]{SS} through the functors
  \cite[(3)]{SS}.  Proposition~\ref{prop:common} supplies the sector on
  which both expansions used in that translation are valid.
\end{proof}

Three boundaries should be stated explicitly.

First, \((r,s)=(1,0)\) is not covered.  By
Remark~\ref{rem:boundary} the formal expression \eqref{eq:ss35} has a
\(z\)-order-zero term there and is not \(J\)-normalized.  This is the
rank-one projective-bundle endpoint, where \(\PP(V)=Z\), \(X'=\varnothing\),
and the point fibres contain no line class \([L]\).  It therefore lies
outside both the normalization argument and the geometric extremal-line
setup.  All genuine discrepancy-one standard flips have \(s\ge1\) and are
covered by Theorem~\ref{thm:assembled}.

Second, everything above lives on the extremal slice.  The Novikov
parameter is restricted to multiples of the class \([L]\) of a line in a
fibre of \(\PP(V)\), all other Novikov variables are zero, and the parameter
is \(\bar{\mathbf t}=t+\ln(q)c_{1}(T)\) with \(t\) pulled back from
\(H^{*}(Z)\); this is the specialization in which \cite[Theorems~1.2
and~1.4]{SS} are stated, and no statement here concerns the unrestricted
quantum product or a comparison across the full Novikov ring.

Third, the two corrections are local to the two hypotheses named in
Section~\ref{sec:hypotheses}.  For \(\nu>1\) nothing changes: the argument
of \cite{SS} applies as printed, with \eqref{eq:ss64} supplied by the
\(\epsilon=1\) branch of \cite[Theorem~A.1]{SS}.  Every other
ingredient of \cite{SS} --- the ring homomorphism to the local model, the
spectral computation in the abstract presentation, the deformation to the
normal cone, the Fourier--Mukai comparison --- is used exactly as it stands.

The count in Lemma~\ref{lem:order} is what makes the discrepancy-one range
behave like the higher-discrepancy one.  The \(s\) numerator factors
\(\sigma_{j}-H\) at \(m=0\), free of \(z\), are worth exactly \(s\) powers,
and at \(\nu=1\) they are the entire margin: the maximal order
\(1-s-\nu d\) sits at \(-1\) precisely when \(s=1\) and \(d=1\), and rises
to \(0\) the moment \(s\) drops to zero.  So the boundary between the
identification \(I=J\) and a nontrivial mirror map at \(\nu=1\) is not the
codimension of the centre but the presence of at least one factor in
\(V'\).
